\section{Limitations}
\label{app:limitations}

\textbf{Scale.} Of 111 designed cases,
90 are retained (Appendix \ref{app:construction}), with
1343 single-primitive instances per model and
defense. Per-target estimates rest on between
159 and
262 instances. These counts are not
directly comparable with benchmarks that cross independent tasks and
attacks: AgentDojo crosses 97 user tasks with 27 injection tasks into 629
security test cases \citep{debenedetti2024agentdojo}, and InjecAgent crosses 17
user cases with 62 attacker cases into 1,054 \citep{zhan2024injecagent}. Each of our
cases pairs one task with an attack designed for it.

\textbf{Domains.} The three environments
are simulated business settings: productivity, DevOps, and a marketplace.
Other benchmarks cover web and computer-use agents
\citep{evtimov2025wasp, kuntz2025osharm}; we do not claim that our
results transfer to them.

\textbf{Attacker.} The attacker of Section \ref{sec:threat} plants content once,
before the run, knowing the environment's public structure but not the user's
task, the system prompt, or records outside its role. Attackers who adapt
content to a deployed defense or to the agent's responses, or who interact
over several turns, are out of scope. The only optimization we apply rewrites
instances against a carrier-only probe (Appendix \ref{app:detection}).
Attacks alter or add calls to tools the task requires.

\textbf{Defense adaptations and measurement.} The runtime defenses are our
adaptations of published systems to a single agent loop, so their results describe these
adaptations and not necessarily the original systems. Defense architectures
we did not evaluate, such as FIDES \citep{costa2025securing} and
AuthGraph \citep{authgraph2026}, are out of scope.
Each instance runs once per configuration at temperature 0, and the
intervals of Appendix \ref{app:uncertainty} resample cases, not runs.

\textbf{Coverage of secondary analyses.} The headline ASR and TSR cover all
five models. Response annotation covers two
models and the first instance of each case--target pair
(Appendix \ref{app:annotation}), the runtime-mechanism logs two
models (Appendix \ref{app:runtime}), and the detection analysis
four (Appendix \ref{app:detection}).

\textbf{Construction with an LLM.} Cases and variants were drafted with Claude
Code under author-specified procedures (Appendix \ref{app:construction}), and two evaluated agents and the
annotation judge are Claude models. We cannot rule out that this shapes the
results. Without a defense, Claude Sonnet 4.5 has the lowest ASR of the five
models (11.1\%), and Claude Sonnet 4 the second
highest (20.6\%), so the Claude models do not share
one direction.

\textbf{Functional, not mechanistic.} The functional targets describe the role a
belief plays in a decision. Models need not represent these roles separately,
and the framework is meant to predict defense coverage, not to describe
internal representations.
